\subsection{The Topological Big Bang: A Phase Transition in Imaginary Time}
Based on the topological analysis (following Perelman and Wick rotation arguments), the "Big Bang" is reinterpreted not as a singularity in real time, but as a \textbf{Topological Crystallization} event.

\subsubsection{Mechanism: The "Big Condensation"}
\begin{enumerate}
    \item \textbf{Pre-Geometric Phase (Layer 3):} The universe begins as a "Quantum Foam" of unconnected logical qubits (random phase, high symmetry $G_{full}$). Topological defects (vortices) are abundant and fluctuating.
    \item \textbf{Crystallization (Layer 4):} As the "Logic Temperature" drops below the critical threshold $T_c$, the superfluid condensate forms. The metric emerges ($g_{\mu\nu} \sim \eta_{\mu\nu}$).
    \item \textbf{Trapped Topology (Origin of Matter):} Crucially, not all defects annihilate during this cooling. A residual density of topological defects is "frozen" into the spacetime manifold due to the non-trivial topology of the vacuum ($\pi_1(T^2) \neq 0$).
    \item \textbf{Simulation Proof:} Our lattice simulation of this Kibble-Zurek mechanism confirms that for a $64 \times 64$ lattice, approximately 34\% of initial defects remain trapped as stable relics (particles) after the transition ($N_i \approx 1320 \to N_f \approx 450$).
\end{enumerate}

\subsubsection{Resolution of the Singularity}
In imaginary time $\tau = it$, the geometry at $t=0$ is a smooth "cap" (like the pole of a sphere), resolved by the Ricci Flow of the condensate. The singularity is an artifact of forcing a Lorentzian description onto a Euclidean phase transition.
